%! @@TITLE@@ — Unit @@UNITINT@@, Lesson @@LESSONID@@ — Student Packet Cover Sheet
\documentclass[12pt]{article}
\usepackage{@@PREFIX@@-article}
\usepackage{@@PREFIX@@-boxes}
\usepackage{ltablex}
\keepXColumns
% Measured banner: the band grows to fit the title block, so a title long enough to wrap grows
% the band rather than running out of it (the stats course's \coverbanner; saar-article does not
% carry one, so it is defined here — every cover copies this block verbatim).
\newsavebox{\bannerbox}
\newlength{\bannerht}
\newlength{\coverbannerpad}
\setlength{\coverbannerpad}{0.16in}       % breathing room above and below the text
\newcommand{\coverbanner}[2]{%
  \sbox{\bannerbox}{%
    \begin{minipage}{\dimexpr\paperwidth-1.4in\relax}%
      \centering
      {\color{white}\textbf{\LARGE \CourseName}}\\[4pt]
      {\color{frostmid}\large\textbf{#1}}\\[6pt]
      {\color{white}\textbf{\large #2}}%
    \end{minipage}}%
  \setlength{\bannerht}{\dimexpr\ht\bannerbox+\dp\bannerbox+2\coverbannerpad\relax}%
  \noindent\begin{tikzpicture}[remember picture, overlay]
    \fill[cerulean] (current page.north west)
      rectangle ([yshift=-\bannerht]current page.north east);
    \node[anchor=north, inner sep=0pt]
      at ([yshift=-\coverbannerpad]current page.north) {\usebox{\bannerbox}};
  \end{tikzpicture}\par
  % The text body starts 0.6in below the page edge (geometry top=0.6in), so reserve only what
  % the band overruns that, plus a gap under the band.
  \vspace*{\dimexpr\bannerht-0.6in+0.16in\relax}%
}

\begin{document}

% Full-bleed cerulean banner, sized to the title block.
\coverbanner{Unit @@UNITINT@@}{Lesson @@LESSONID@@ \quad @@TITLE@@}

\vspace{0.06in}
% NAMESTRIP: this is the ONLY component that carries the name row. Everything else
% is stapled behind it.
\namedateperiod
\vspace{0.18in}

% THERE IS NO SPOILER RULE. The vocabulary is taught, so name it outright here — "standard
% deviation", not "how spread out the values are". One target per priority idea/skill, worded
% exactly as the notes' objectivebox words them. The lettered standard codes belong in the
% teacher-facing lesson plan, not here.
\begin{learningtargetbox}
\small
\begin{enumerate}[leftmargin=1.5em, itemsep=5pt]
  \item \ldots{} TODO: learning target 1, naming the formal vocabulary.
  \item \ldots{} TODO: learning target 2, naming the formal vocabulary.
\end{enumerate}
\end{learningtargetbox}

\vspace{0.14in}

\begin{tocbox}
\small
\renewcommand{\arraystretch}{1.5}
\begin{tabularx}{\linewidth}{@{} c l X r @{}}
\rowcolor{frost}
\textbf{\#} & \textbf{Component} & \textbf{Description} & \textbf{Score} \\
\midrule
% THREE ROWS. The Guided Practice has no handout of its own — it is inside the notes — so it
% does not get a row; it is named inside the Guided Notes & Practice description.
1 & Warm-Up    & TODO --- what the spiral items rehearse & \blank{1.2cm} \\
2 & Guided Notes \& Practice & TODO --- the ideas the two sections build, then the one survey we
                 work together & \blank{1.2cm} \\
% Row 3 is the lesson's GRADED practice — the individual practice, started in class — so the
% score cell carries a \blank{} — never NA. Whether it is assigned from this packet or swapped
% for an equivalent DESMOS activity is decided lesson by lesson at Close & Assign; the pages are
% authored either way, so the row and its score cell never change.
3 & Homework   & TODO --- the new context in one line --- \textbf{scored, started in class, due the
                 first class after two study halls} & \blank{1.2cm} \\
\midrule
  & \multicolumn{2}{r}{\textbf{Total}} & \blank{1.2cm} \\
\end{tabularx}
\end{tocbox}

\vspace{0.14in}

% Keep in Mind is the lesson's CONTENT takeaway, in the formal vocabulary: the rule, the test
% to apply, and the trap to avoid. This is the page a student flips back to while doing the
% homework, so make it usable on its own.
\begin{remindbox}
\small \textbf{TODO: the rule, in one sentence, using the formal term.} TODO: the test to apply
--- how you decide which case you are in. \textbf{Watch out:} TODO --- the trap, and what it
looks like when you have fallen into it.
\end{remindbox}

\end{document}
